\subsection{Redes neuronales convolucionales}

Las redes neuronales convolucionales (\acrshort{cnn}) constituyen una de las arquitecturas de aprendizaje profundo más empleadas para el análisis de datos complejos. Aunque fueron desarrolladas originalmente para tareas de visión artificial, su capacidad para extraer características relevantes de manera automática y eficiente las ha convertido en una herramienta ampliamente utilizada en ámbitos como: el procesamiento de imágenes, el reconocimiento de voz, el procesamiento del lenguaje natural y, más recientemente, el análisis de señales biomédicas como el electroencefalograma (\acrshort{eeg}) \cite{albawi_UnderstandingConvolutional2017,lecun_DeepLearning2015,edelman_NonInvasiveBrainComputer2025,schmidhuber_DeepLearning2015}.

La operación en la que se basan estas redes es la convolución, una transformación matemática que aplica pequeños filtros o \textit{kernels} sobre la entrada para detectar patrones específicos. Cada filtro aprende durante el entrenamiento a responder a determinadas características presentes en los datos, generando como salida un conjunto de mapas de características (\textit{feature maps}). Estos \textit{feature maps} se vuelven progresivamente más abstractos y complejos conforme se avanza en la red, lo que le permite aprender automáticamente las características más discriminativas para la tarea de clasificación. Las capas convolucionales requieren un número significativamente menor de parámetros que las capas completamente conectadas (\textit{fully connected}), reduciendo el coste computacional y favoreciendo la capacidad de generalización del modelo \cite{albawi_UnderstandingConvolutional2017,edelman_NonInvasiveBrainComputer2025,lecun_DeepLearning2015,goodfellow_DeepLearning2016}.

Habitualmente, una CNN está formada por una sucesión de capas convolucionales y capas de reducción dimensional (\textit{pooling}), seguidas de una o varias capas \textit{fully connected}. Las capas convolucionales se encargan de extraer características relevantes de la señal de entrada, mientras que las capas de \textit{pooling} reducen la dimensionalidad de las representaciones obtenidas y contribuyen a disminuir el sobreajuste. Finalmente, las capas \textit{fully connected} utilizan las características extraídas para realizar la clasificación o regresión correspondiente \cite{albawi_UnderstandingConvolutional2017,rakhmatulin_ExploringConvolutional2024}.

En el ámbito de las interfaces cerebro-computador (\acrshort{bci}), las CNN se han consolidado como una de las técnicas más utilizadas para la decodificación de señales EEG \cite{craik_DeepLearning2019,schirrmeister_DeepLearning2017,edelman_NonInvasiveBrainComputer2025}. A diferencia de las imágenes, las señales EEG presentan simultáneamente una dimensión temporal y una dimensión espacial asociada a la distribución de los electrodos sobre el cuero cabelludo. Por este motivo, las arquitecturas diseñadas para EEG suelen combinar convoluciones temporales, encargadas de extraer información espectral y patrones dependientes del tiempo, con convoluciones espaciales que modelan las relaciones entre los distintos canales de la señal \cite{edelman_NonInvasiveBrainComputer2025}.

Tradicionalmente, los sistemas BCI basados en EEG dependían de técnicas de extracción manual de características combinadas con clasificadores convencionales. Sin embargo, la aparición de arquitecturas de redes neuronales profundas específicas para EEG, como DeepConvNet y ShallowConvNet, supuso un cambio de paradigma al integrar de forma conjunta las etapas de extracción de características y clasificación dentro de una misma red neuronal \cite{schirrmeister_DeepLearning2017,kollod_DeepComparisons2023}.
